% Source: https://artofproblemsolving.com/community/c3h2593161_complex_number_multiplication_on_a_complex_plane
% Stable identity: aops-topic-2593161-q01
\textbf{(a)} Multiply out, using $i^2=-1$:
\[wz=rs\bigl[(\cos\theta\cos\varphi-\sin\theta\sin\varphi)+i(\sin\theta\cos\varphi+\cos\theta\sin\varphi)\bigr]=rs\bigl(\cos(\theta+\varphi)+i\sin(\theta+\varphi)\bigr)\]
by the angle-addition formulas.

\textbf{(b)} By (a), the point $z$ at distance $s$ from $0$ and angle $\varphi$ goes to the point at distance $rs$ and angle $\varphi+\theta$. Rotating $z$ about $0$ through $\theta$ gives the point at distance $s$ and angle $\varphi+\theta$; dilating by factor $r$ from $0$ then gives distance $rs$ at the same angle. So $z\mapsto wz$ is exactly this rotation followed by this dilation. (Both maps send $0$ to $0$.) In short: lengths multiply and angles add. Addition, by contrast, is a translation.

\textbf{(c)} $w=1+i=\sqrt2\left(\cos45^\circ+i\sin45^\circ\right)$, so multiplying by $w$ rotates by $45^\circ$ and scales by $\sqrt2$. Indeed
\[(1+i)(2+i)=2+i+2i+i^2=\boxed{1+3i},\]
and $|1+3i|=\sqrt{10}=\sqrt2\cdot\sqrt5=|w|\,|z|$. $\blacksquare$

\begin{asy}
usepackage("amsmath");
size(4.5cm);
defaultpen(fontsize(10pt));
pair O=(0,0),z=(2,1),w=(1,3);
draw((-0.6,0)--(3,0),gray,Arrow);draw((0,-0.4)--(0,3.4),gray,Arrow);
draw(O--z,linewidth(0.8),Arrow);draw(O--w,linewidth(0.8),Arrow);
draw(z--w,dashed);
draw(arc(O,0.9,degrees(z),degrees(w)),Arrow);
label("$45^\circ$",0.9*dir((degrees(z)+degrees(w))/2),NE,fontsize(8pt));
dot(z);dot(w);
label("$z=2+i$",z,E);label("$(1+i)z=1+3i$",w,E);
label("$\sqrt5$",z/2,SE);label("$\sqrt{10}$",w/2,W);
label("$\mathrm{Re}$",(3,0),E);label("$\mathrm{Im}$",(0,3.4),W);
\end{asy}
